% year: תשפ"ג
% semester: א
% moed: ג
% source: exams/מבחנים/תשפג/מועד ג תשפג.pdf
% number: 2א
% categories: taylor

\begin{question}
השתמשו בקירוב לינארי כדי להעריך את הביטוי $\log_2 10$.
\end{question}

\begin{hint}
בחרו נקודה קרובה ל-$10$ שבה $\log_2$ ידוע, למשל $x_0=8$.
\end{hint}

\begin{solution}
נגדיר $f(x)=\log_2 x=\frac{\ln x}{\ln 2}$, ואז $f'(x)=\frac{1}{x\ln 2}$. הקירוב הלינארי סביב $x_0$ הוא
\[ f(x)\approx f(x_0)+f'(x_0)(x-x_0). \]
ניקח $x_0=8$ (שם $\log_2 8=3$) ו-$x=10$:
\[ \log_2 10\approx 3+\frac{1}{8\ln 2}\cdot 2=3+\frac{1}{4\ln 2}\approx 3+0.3607=\boxed{3.3607}. \]
(הערך המדויק $\log_2 10\approx 3.3219$; הקירוב גבוה מעט מהערך האמיתי, כצפוי, כי $\log_2$ קעורה ולכן המשיק נמצא מעל הגרף.)
\end{solution}
